% Part: proof-theory
% Chapter: cut-elimination
% Section: introduction

\documentclass[../../../include/open-logic-section]{subfiles}

\begin{document}

\olfileid{pt}{cut}{int}

\olsection{ভূমিকা}

\Cut{} বিধিটি হলো
\[
\Axiom$\Gamma \fCenter \Delta, !A$
\Axiom$!A, \Pi \fCenter \Lambda$
\RightLabel{\Cut}
\BinaryInf$\Gamma, \Pi \fCenter \Delta, \Lambda$
\DisplayProof
\]
(!!{formula}~$!A$-কে \emph{কর্তনসূত্র} বলে।)

জেন্টজেনের মূল সিকোয়েন্ট কলন~\Log{LK}-তে \Cut{} ছিল।
জেন্টজেনের \emph{Hauptsatz} (``মূল উপপাদ্য'') বলে,
\Log{LK}-তে !!{provable} প্রতিটি সিকোয়েন্ট \Cut{} বিধি
ছাড়াও !!{provable}; অর্থাৎ \Log{LK}-প্রমাণ থেকে
\Cut{} ``বর্জন'' করা যায়। আমাদের \Log{G1c} ও \Log{G3c}
পদ্ধতিতে \Cut{} বিধি নেই। তবু সেগুলির ক্ষেত্রেও প্রশ্ন করা যায়:
\Log{G1c} (বা \Log{G3c})-এর অন্য বিধির সঙ্গে \Cut{} ব্যবহার করা
!!{proof}s থেকে কি \Cut{} বর্জন করা যায়? প্রতিষ্ঠিত পরিভাষায়
সমতুল্য প্রশ্নটি হলো: \Log{G1c} (বা \Log{G3c})-তে \Cut{}
কি গ্রহণযোগ্য বিধি?

কর্তন-বর্জন উপপাদ্য প্রমাণতত্ত্বের সম্ভবত সবচেয়ে গুরুত্বপূর্ণ
ও সুপরিচিত ফল। প্রথমত, এটি দেখায় যে প্রথম দর্শনে আবশ্যক মনে
হলেও \Log{G1c} ও \Log{G3c}-তে একটি বিধি লাগে না। বাস্তব
প্রমাণ আনুষ্ঠানিক করতে গেলে লেমা ব্যবহারের ব্যবস্থা দরকার
হয়। গণিতবিদেরা একটি ফল প্রমাণ করে পরে অন্য ফলের প্রমাণে
সেটি ব্যবহার করেন। ধরি, প্রথমে লেমা~$L$-এর প্রমাণকে
$\Sequent L$ সিকোয়েন্টের !!{proof} হিসেবে আনুষ্ঠানিক করলাম।
পরে $L$ ব্যবহার করে ফল~$R$-এর প্রমাণকে $L \Sequent R$
সিকোয়েন্টের !!{proof} হিসেবে আনুষ্ঠানিক করলাম। কেবল
$\Sequent R$-এর, অর্থাৎ~$R$-এর, !!a{proof} আছে কীভাবে জানব?
দুই !!{proof}s মিলিয়ে দেওয়ার ব্যবস্থা চাই; \Cut{} বিধি তা-ই দেয়।

আগে বলেছি, \Log{G1c} ও \Log{G3c} পূর্ণ; তাই প্রত্যাশিত
সব প্রমাণযোগ্য সিকোয়েন্ট—অর্থাৎ সব বৈধ সিকোয়েন্ট—তারা
প্রমাণ করে। এটি সরাসরিও প্রমাণ করা যায়। কিন্তু জেন্টজেনের
সময়ে শুধু স্বতঃসিদ্ধমূলক নিষ্পাদন-ব্যবস্থার পূর্ণতা জানা ছিল।
\Log{LK}-র পূর্ণতা দেখাতে জেন্টজেন স্বতঃসিদ্ধমূলক ব্যবস্থার
!!{derivation}s-কে \Log{LK}-!!{proof}s-এ রূপান্তর করেন।
সেই রূপান্তরে \Cut{} অপরিহার্যভাবে ব্যবহৃত হয়।
কর্তন-বর্জন শুধু ধ্রুপদি যুক্তিবিদ্যাতেই গুরুত্বপূর্ণ নয়:
অন্য অনেক যুক্তিবিদ্যারও সিকোয়েন্ট কলন আছে। কোনো কোনো
যুক্তিবিদ্যায় সিকোয়েন্ট কলনের পূর্ণতা সরাসরি প্রমাণ করা কঠিন
বা অসম্ভব; তখন অন্য ব্যবস্থা থেকে \Cut{} ব্যবহার করে
রূপান্তরই পূর্ণতা দেখানোর একমাত্র পথ। পরে প্রমাণতাত্ত্বিকভাবে
কর্তন-বর্জন দেখানো দরকার, যাতে বোঝা যায় \Cut{} অনাবশ্যক
এবং \Cut{}-বিহীন ব্যবস্থাও পূর্ণ।

কর্তন-বর্জন থেকে স্বতন্ত্রভাবে আকর্ষণীয় আরও ফল পাওয়া যায়।
এখানে দুটি উদাহরণ দেখব: ক্রেগের ইন্টারপোলেশন লেমা ও
Herbrand-এর উপপাদ্য।

কর্তন-বর্জনের প্রমাণে সাধারণত \Cut{}-যুক্ত !!a{proof}-কে
তা-বিহীন প্রমাণে রূপান্তরের উপায় দেখানো হয়। এই রূপান্তরগুলি
সাধারণত ``স্থানীয়'': !!a{proof}-এর কোনো অংশ নিই—সাধারণত
\Cut{} অনুমানে শেষ হওয়া উপ-!!{proof}—এবং সেটিকে একই সিকোয়েন্টের
আরেকটি উপ-!!{proof} দিয়ে বদলাই, যেখানে \Cut{} অনুমানটি
সরানো হয়েছে বা অন্য অনুমান দিয়ে প্রতিস্থাপিত। এই যুক্তি
\Cut{}-যুক্ত !!{proof}s-কে তা-বিহীন প্রমাণে রূপান্তরের ধাপে ধাপে
অ্যালগরিদম দেয়। প্রয়োগে তা আরও তথ্য দেয়। যেমন,
ইন্টারপোলেশন লেমার মডেলতাত্ত্বিক প্রমাণে বলা যেতে পারে:
``$!A \lif !B$ বৈধ হলে একটি ইন্টারপোল্যান্ট আছে''
(আপাতত ইন্টারপোল্যান্ট কী, তা থাক)। বিপরীতে, কর্তন-বর্জন
ব্যবহারকারী প্রমাণতাত্ত্বিক যুক্তি $!A \lif !B$-র !!a{proof}
থেকে ইন্টারপোল্যান্ট গণনার অ্যালগরিদম দেয়। এই অর্থে
প্রমাণতাত্ত্বিক যুক্তিটি \emph{নির্মাণমূলক}।
% BN-SRC-530 TARGET-CORRECTION-BEGIN
\par\smallskip\noindent\textnormal{\small\textbf{উৎস-সংশোধন (BN-SRC-530):} উৎসে একটি প্রমাণ-বিশেষক টোকেন দুই লাইনে ভেঙে গিয়েছিল; লক্ষ্যপাঠে সেটি অখণ্ড করা হয়েছে।}
% BN-SRC-530 TARGET-CORRECTION-END

কর্তন-বর্জন উপপাদ্য প্রমাণের দুটি প্রচলিত পথ আছে। একটি
(জেন্টজেন থেকে আসা) প্রমাণের সর্বোচ্চে-থাকা কর্তন
সরায়, তারপর আর কর্তন না-থাকা পর্যন্ত এমন কর্তন সরিয়ে যায়।
অন্যটি (Sch\"utte ও Tait-এর প্রবর্তিত) !!a{proof}-এর সবচেয়ে
জটিল কর্তনগুলিতে নজর দেয় এবং পর্যায়ক্রমে এমন কর্তন সরায়
যেগুলির কর্তন !!{formula}-র গভীরতার চেয়ে বেশি গভীরতার কর্তনসূত্র
আর কোনো কর্তনে নেই। উভয় পদ্ধতির সুবিধা ও অসুবিধা আছে।
কোনো ব্যবস্থায় একটি পথ চলে, অন্যটি কঠিন বা অসম্ভব হতে
পারে। তাই দুটিই বর্ণনা করব। \Log{G3c}-তে কর্তন-বর্জন
প্রমাণ করব। প্রকৃতপক্ষে সাধারণ \Cut{} নয়,
\emph{অভিন্ন-প্রসঙ্গ কর্তন}~\CutCS{} অপসারণ দেখাব:
\[
\Axiom$\Gamma \fCenter \Delta, !A$
\Axiom$!A, \Gamma \fCenter \Delta$
\RightLabel{\CutCS}
\BinaryInf$\Gamma \fCenter \Delta$
\DisplayProof
\]
% BN-SRC-529 TARGET-CORRECTION-BEGIN
\par\smallskip\noindent\textnormal{\small\textbf{উৎস-সংশোধন (BN-SRC-529):} অভিন্ন-প্রসঙ্গ কর্তনের প্রদর্শিত বিধিতে ভুল সাধারণ কর্তন-লেবেলের বদলে অভিন্ন-প্রসঙ্গ কর্তনের লেবেল বসানো হয়েছে।}
% BN-SRC-529 TARGET-CORRECTION-END
কোনো সিকোয়েন্ট কলনে দুর্বলীকরণ ও সংকোচন থাকলে
(\Log{G1c}-র মতো প্রকৃত বিধি হিসেবে অথবা \Log{G3c}-র মতো
গ্রহণযোগ্য বিধি হিসেবে), \Cut{} দিয়ে \CutCS{}-কে, আবার
\CutCS{} দিয়ে \Cut{}-কে অনুকরণ করা যায়। \Cut-এর বদলে
\CutCS{} বেছে নিচ্ছি, কারণ প্রথম কর্তন-বর্জন কৌশলটি
\CutCS{}-এর জন্য বর্ণনা করা সহজ, আর দ্বিতীয় কৌশলটি
শুধু \CutCS{}-এর জন্যই চলে।

\end{document}
